
\section{Round-fourteen positive closure: the joint collision Gaussian form}
\label{sec:r14-b3}

The Gaussian field is constructed first from finite-volume joint cumulants.
It is not inferred from a formal Hessian of the action.  Its collision driver
acts on the complete coefficient \(\Delta p+\psi\), so pure contact
observables have nonzero variance.  Tightness is proved from unconditional
deterministic-interval estimates and a multiscale grid; no stopping-time
estimate conditioned on the complete deterministic microstate is used.

\subsection{The collision quotient and balanced source space}

Let
\[
 \mathfrak C_T=
 [0,T]\times\mathbb T^3\times\mathbb R^3\times\mathbb R^3
 \times\mathbb S^2/\!\sim ,
\qquad
 (v,v_*,\omega)\sim(v_*,v,-\omega),
\]
and also identify pre/post variables by the elastic involution when a test is
declared symmetric.  Let \(\mathcal Y_\psi\) be the exponential Orlicz heart
for the reference contact measure \(A_f\).  A density test \(p\) belongs to
\(\mathcal Y_p\) when its endpoint traces, transport derivative, and
\[
 \Delta_\omega p
 =p(v')+p(v_*')-p(v)-p(v_*)
\]
belong to the corresponding weighted spaces.

Define
\[
 \mathfrak D(p,\psi)=\Delta p+\psi.
\]
The full finite-horizon balance coboundary generated by \(r\) is
\[
 \mathfrak G r=
 \left(
 r,\,-\Delta r,\,
 -r(0),\,r(T),\,
 -(\partial_t+v\cdot\nabla_x)r
 \right),
\]
with all endpoint and transport components included.

\begin{theorem}[Closed balanced gauge]
\label{thm:r14-b3-gauge}
The weak balance operator from density/contact paths to endpoint, transport,
and collision distributions has closed range on the short-time weighted
Maxwellian chart.  Its annihilator is exactly the closure of the full
coboundaries \(\mathfrak G r\) together with the five global conservation
constraints.  On a balanced pair the dual action depends only on
\(\mathfrak D(p,\psi)\).
\end{theorem}

\begin{proof}
Solve the backward transport equation with terminal data on the torus and
use the collision Dirichlet form on each velocity fiber.  On the complement
of the five invariants, the Maxwellian lower bound and the small bounded
collision tilt give
\[
 \int|\Delta r|^2\,dA_f
 +\|(\partial_t+v\cdot\nabla_x)r\|_{-1,w}^2
 \ge c\|r\|_{\mathcal Y_p/\mathcal I}^2 .
\]
This is a finite-time observability estimate, not a long-time Liv\v{s}ic
claim.  The closed-range theorem gives the asserted range and annihilator.
Substitution of the weak balance identity reduces every dual pair to the
single collision coefficient \(\Delta p+\psi\).
\end{proof}

\subsection{The regular biased chart}

Let \(q=e^{\Delta p+\psi}\), and let
\[
 dA_f={1\over2}ff_*((v-v_*)\cdot\omega)_+
 dt\,dx\,dv\,dv_*\,d\omega .
\]

\begin{lemma}[Perturbative nonautonomous collision estimate]
\label{lem:r14-b3-energy}
On a sufficiently small bounded source chart produced by B2, the density
satisfies
\[
 cM(v)\le f(t,x,v)\le CM(v),\qquad
 0<q_-\le q,q^*\le q_+,
\]
where \(q^*\) is the pre/post transform.  The symmetric part of the
linearized collision operator has Dirichlet form comparable with
\[
 {1\over4}\int
 (q+q^*)ff_*
 \bigl(h'+h_*'-h-h_*\bigr)^2
 ((v-v_*)\cdot\omega)_+ ,
\]
and is coercive modulo the five velocity collision invariants.  The
nonautonomous transport--collision equation is well posed on
\([0,T_0]\) and obeys a weighted finite-time energy estimate.
\end{lemma}

\begin{proof}
B2 constructs the source-biased solution by Picard iteration from a
Maxwellian-comparable datum.  Small time and a bounded source preserve the
two-sided Maxwellian comparison and the positive bounds on \(q,q^*\).
Average the collision bilinear form with its pre/post adjoint.  The displayed
nonnegative square is the leading term; all coefficient differences from the
equilibrium form are relatively bounded by the source radius.  The classical
Maxwellian collision gap is therefore stable under this perturbation.

For spatially dependent coefficients, do not diagonalize by Fourier
transform.  Commute the equation with the hypocoercive multiplier
\(\nabla_x(1-\Delta_x)^{-1}\nabla_v\), estimate coefficient commutators by
their bounded chart norms, and integrate the resulting differential
inequality.  The transport part is skew after the commutator correction, the
microscopic part is controlled by the collision gap, and the five local
moments satisfy the exact finite hyperbolic balance coupled to the controlled
higher moments.  Gronwall gives the finite-time estimate.
\end{proof}

\subsection{Covariance from microscopic cumulants}

Under a regular biased microscopic law, define
\[
 Z_\varepsilon(p,\psi)
 =
 \sqrt{\mu_\varepsilon}
 \left[
 \langle p,\pi^\varepsilon-f\rangle
 +\langle\psi,\Gamma^\varepsilon-\Gamma\rangle
 \right].
\]
B2 gives normal convergence of the joint source pressure on a complex
neighborhood.

\begin{theorem}[Joint covariance kernel]
\label{thm:r14-b3-covariance}
The finite-dimensional laws of \(Z_\varepsilon\) converge to a centered
Gaussian field \(Z\).  Its dynamic martingale bracket is
\[
 \begin{aligned}
 &\langle M(p,\psi),M(\tilde p,\tilde\psi)\rangle_t\\
 &\quad=
 \int_0^t\int
 q\,(\Delta p+\psi)(\Delta\tilde p+\tilde\psi)\,dA_f ,
 \end{aligned}
\]
and its initial covariance is the B1 microcanonical Schur complement.
Equivalently, there is an isonormal Gaussian random measure
\(\mathbb W\) on \(L^2(qA_f)\) such that the direct collision noise is
\[
 M_t(p,\psi)=
 \int_{[0,t]\times\mathfrak C}
 (\Delta p+\psi)\sqrt q\,d\mathbb W .
\]
In particular a pure contact test has variance
\(\int q\psi^2\,dA_f\).
\end{theorem}

\begin{proof}
For \(k\ge3\), the \(k\)-th centered cumulant of the
\(\sqrt{\mu_\varepsilon}\)-scaled field is
\(O(\mu_\varepsilon^{1-k/2})\), by the normal B2 connected expansion.
The second cumulant converges by two source derivatives.  Differentiating the
collision exponential
\(e^{\Delta p+\psi}\) twice gives exactly the displayed product, including
contact--contact and density--contact terms.  B1 supplies the conditioned
initial block.  The limiting positive bilinear form defines the isonormal
measure and the Gaussian finite-dimensional laws.
\end{proof}

The linear density fluctuation is the unique solution of the nonautonomous
linearized biased Boltzmann equation driven by this collision measure.
Lemma~\ref{lem:r14-b3-energy} gives its continuous solution operator.

\subsection{Deterministic-interval tightness}

Let \(\mathcal F_t^{\rm obs}\) be the filtration generated by the empirical
density/contact fields, not the complete microscopic state.  We do not claim
a Poisson conditional intensity in this filtration.

\begin{lemma}[Localized cumulant moments]
\label{lem:r14-b3-moments}
For every smooth compact family of tests and every deterministic interval
\(I\) of length \(h\),
\[
 \mathbb E|Z_\varepsilon(I)|^{2m}
 \le C_m\left(h^m+
 \sum_{\ell=1}^{m-1}\mu_\varepsilon^{-\ell}h^{m-\ell}\right).
\]
In particular
\[
 \mathbb E|Z_\varepsilon(I)|^4
 \le C\left(h^2+{h\over\mu_\varepsilon}\right).
\]
The estimate is unconditional and remains valid jointly for density and
contact increments.
\end{lemma}

\begin{proof}
Insert time-supported sources into the B2 connected expansion.  A connected
cumulant with \(j\) distinct collision times has an \(L^1\) density on
\(I^j\), hence contributes \(O(h^j)\).  Coincident contact insertions produce
diagonal contractions and one lost factor of \(\mu_\varepsilon\).  Summing
partitions of \(2m\) insertions gives the displayed powers.  This is the
correct microscopic diagonal term; no estimate conditioned just before a
deterministically known collision is asserted.
\end{proof}

\begin{theorem}[Process-level joint Gaussian limit]
\label{thm:r14-b3-process}
The density/contact fluctuation pair is \(C\)-tight in the weighted
distribution path space and converges to the Gaussian solution determined by
Theorem~\ref{thm:r14-b3-covariance}.
\end{theorem}

\begin{proof}
Take a dyadic grid with mesh \(\delta\).  Markov's inequality and
Lemma~\ref{lem:r14-b3-moments} with \(m\ge3\) give
\[
 \mathbb P\left(
 \max_{I\in\mathcal D_\delta}|Z_\varepsilon(I)|>\eta
 \right)
 \le C_\eta(\delta^{m-1}+\mu_\varepsilon^{-1}\delta^{m-2}+\cdots).
\]
Sum over dyadic levels down to \(\mu_\varepsilon^{-2}\).  Below that scale,
the maximal jump is \(O(\mu_\varepsilon^{-1/2})\), while the B2 contact-count
exponential moment controls the probability of multiple jumps in one finest
cell.  Letting first \(\varepsilon\to0\) and then \(\delta\to0\) proves
the modulus estimate.  Mitoma's criterion and weighted compact containment
give \(C\)-tightness.  Finite-dimensional convergence and uniqueness of the
linear Gaussian equation identify the limit.
\end{proof}

\subsection{The Cameron--Martin form and the action tangent}

Let \(\Sigma\) be the covariance operator on the quotient dual.  Define the
closed extended form
\[
 \mathfrak q(z)=
 \begin{cases}
 \|\Sigma^{-1/2}z\|^2,
 &z\in\operatorname{Ran}\Sigma^{1/2},\\
 +\infty,&\text{otherwise}.
 \end{cases}
\]

\begin{theorem}[Second epi-derivative]
\label{thm:r14-b3-epi}
At every regular exposed path of the B2 good rate,
\[
 d_e^2I(z)=\mathfrak q(z).
\]
The dual kernel of \(\Sigma\) is the closed full balance-coboundary space of
Theorem~\ref{thm:r14-b3-gauge}; after quotienting it,
\(\mathfrak q\) has no nonzero primal null vector.
\end{theorem}

\begin{proof}
On every finite family of sources, B2 convex duality and the analytic
pressure expansion give the inverse of the finite covariance matrix.
Projective monotonicity follows because enlarging the source family takes a
supremum of finite quadratic duals.  The supremum is precisely the closed
form \(\|\Sigma^{-1/2}z\|^2\) on
\(\operatorname{Ran}\Sigma^{1/2}\).  B2's full LDP supplies the
finite-rate recovery needed for Mosco limsup.  The finite-horizon
observability theorem identifies the dual kernel with exact balance
coboundaries, rather than invoking a long-time Liv\v{s}ic theorem.
\end{proof}

\begin{theorem}[Round-fourteen B3 closure]
\label{thm:r14-b3-main}
The hard-sphere density/contact fluctuations have a correctly typed joint
collision-space Gaussian driver, unconditional deterministic-interval process
tightness, and the Cameron--Martin action tangent on
\(\operatorname{Ran}\Sigma^{1/2}\).
\end{theorem}

\begin{proof}
Combine Theorems~\ref{thm:r14-b3-covariance},
\ref{thm:r14-b3-process}, and \ref{thm:r14-b3-epi}, with
Theorem~\ref{thm:r14-b3-gauge}.
\end{proof}
